\chapter{Supplementary Mathematical Background}
\label{app:mathematical_background}

This appendix collects standard results that support the main text but are not needed to follow its central argument~\cite{nielsen_chuang_2010_quantum_computation_and_quantum_information,quantum_theory_from_first_principles}.  The notation follows Chapter~\ref{ch:MathBackground}.

\section{Generalized Measurements}

\begin{definition}[Kraus operators]
    An outcome-resolved quantum measurement can be described by Kraus operators $\{K_{x,a}\}$ with
    \begin{equation}
        p(x|\rho)=\sum_a\operatorname{Tr}\!\left[K_{x,a}\rho K_{x,a}^\dagger\right]
        \qquad
        \sum_{x,a}K_{x,a}^\dagger K_{x,a}=\mathbb{I}
    \end{equation}
    The associated POVM elements are $M_x=\sum_aK_{x,a}^\dagger K_{x,a}$, as in Eq.~\eqref{eq:povm_def}.  A POVM element is also called an effect and is often written as $E_x=M_x$.

    The POVM gives only the probability of each outcome.  The Kraus operators also describe the state after that outcome.  If outcome $x$ is observed and $p(x|\rho)>0$, the normalized state is
    \begin{equation}
        \rho_x
        =\frac{\sum_aK_{x,a}\rho K_{x,a}^\dagger}{p(x|\rho)}
    \end{equation}
    The index $a$ allows several unobserved state changes to produce the same recorded outcome $x$.
\end{definition}

\begin{theorem}[Naimark dilation]
    Every POVM can be implemented as a projective measurement on a larger Hilbert space.  There are an ancilla state $|0\rangle_A$, a unitary $U$, and projectors $\{\Pi_x\}$ such that
    \begin{equation}
        \operatorname{Tr}[\rho M_x]
        =\operatorname{Tr}\!\left[
            U(\rho\otimes|0\rangle\langle0|_A)U^\dagger\Pi_x
            \right]
        \label{eq:naimark_dilation}
    \end{equation}
\end{theorem}

\begin{remark}[Operational meaning]
    Naimark dilation shows that a generalized measurement can be built from ordinary unitary evolution and a projective measurement on a larger system.  This is why optimization over POVMs in the definition of QFI still refers to a physical measurement procedure.
\end{remark}

\section{Choi-Matrix Representation of a Quantum Channel}
\label{sec:choi_quantum_channel}

Chapter~\ref{ch:MathBackground} introduces channels through their action on density operators and through Kraus operators.  Here we give the equivalent Choi-matrix description.  A channel and its Choi matrix contain the same information, but they are different mathematical objects.  The Choi matrix packages the action of the whole channel into one operator.  Let $J(\mathcal{E})$ denote the Choi matrix of a channel $\mathcal{E}:\mathcal{L}(\mathcal{X})\rightarrow\mathcal{L}(\mathcal{Y})$.  In the column Choi convention, the action on an input state is~\cite{wood2015graphical}
\begin{equation}
    \mathcal{E}(\rho)
    =\operatorname{Tr}_{\mathcal{X}}\!\left[
        (\rho^{\mathsf{T}}\otimes\mathbb{I}_{\mathcal{Y}})
        J(\mathcal{E})
        \right]
    \label{eq:choi_channel_action}
\end{equation}


\section{Classical Statistical Distance and Metrological Scaling}

\begin{definition}[Wootters statistical distance]
    For nearby probability distributions, the Wootters statistical distance is
    \begin{equation}
        \mathrm{d}s^2=\frac{1}{4}\sum_x\frac{(\mathrm{d}p_x)^2}{p_x}
    \end{equation}
    For a one-parameter family, $\mathrm{d}s^2=\mathcal{F}_{\mathrm{C}}(\theta)\mathrm{d}\theta^2/4$.  The factor $1/p_x$ means that the same absolute probability change carries more weight for a rare outcome than for a common one.
\end{definition}

\begin{remark}[Metrological scaling]
    Fisher information measures how quickly nearby probability distributions separate when the parameter changes.  For $N$ independent probes with fixed information per probe, the Cram\'er--Rao bound gives the standard-quantum-limit scaling $\Delta\theta=\mathcal{O}(N^{-1/2})$.  Under ideal unitary encoding, entangled probes can reach the Heisenberg scaling $\Delta\theta=\mathcal{O}(N^{-1})$.  Noise, preparation cost, and the required measurement determine if this ideal scaling is useful in practice.
\end{remark}

\section{Geometric and Multiparameter Extensions}

\begin{definition}[Quantum geometric tensor]
    For a pure ansatz state, the complex quantum geometric tensor is
    \begin{equation}
        Q_{\mu\nu}
        =\langle\partial_\mu\psi|
        (\mathbb{I}-|\psi\rangle\langle\psi|)
        |\partial_\nu\psi\rangle
        \label{eq:qgt_def}
    \end{equation}
    Its real part is the Fubini--Study metric used by the quantum natural gradient.  Its imaginary part is the Berry curvature.  In simple terms, the real part measures how far the state moves, while the imaginary part records the geometric phase structure of that motion.
\end{definition}

\begin{remark}[Scope of the main text]
    The main text uses only the real metric component.  The Bures metric belongs to a wider family of contractive metrics classified by Petz.  These metrics express the same basic idea: a parameter-independent quantum channel cannot make two nearby states easier to distinguish.  This Thesis uses the Bures/SLD metric because Eq.~\eqref{def:qfim_bures_relation} connects it directly to QFI.
\end{remark}

\begin{remark}[Multiparameter compatibility]
    For several physical parameters, the inverse QFIM gives a local precision benchmark.  However, the best measurement for one parameter may be incompatible with the best measurement for another.  A common sufficient condition for local compatibility is
    \begin{equation}
        \operatorname{Tr}\!\left[\rho[L_i,L_j]\right]=0
        \quad\text{for all } i,j
        \label{eq:weak_commutativity}
    \end{equation}
    When this condition fails, the Holevo bound gives the tighter asymptotic limit for simultaneous estimation.  These refinements are not needed for the single-parameter QFI estimates and variational optimization studied in this Thesis.
\end{remark}

\section{Rank Boundaries}

\begin{remark}[Rank boundaries]
    At a point where the rank of $\rho(\theta)$ changes, the SLD may not be unique outside the support of the state.  Intuitively, directions with zero probability do not define a unique observable response.  The QFI remains well defined when the spectral sums and vectorized inverse of Chapter~\ref{ch:MathBackground} keep only terms with a nonzero denominator, as in Eqs.~\eqref{eq:sld_spectral_elements} and~\eqref{eq:safranek_vectorized}.  This is also the numerically stable convention used in the classical reference calculations.
\end{remark}
